% ============================================================================
% Fichier  : partie2/P02_C06_ids_ips.tex
% Titre    : La Détection d'Intrusion (IDS/IPS)
% Auteur   : MINKA MI NGUIDJOÏ Thierry Emmanuel
% ============================================================================
% Sources : NIST SP 800-94 (IDS/IPS) ; MITRE ATT&CK.
% ============================================================================

\chapter{La Détection d'Intrusion (IDS/IPS)}
\label{chap:idsips}
\chaptermark{La détection d'intrusion}

\epigraph{%
    « Détecter laisse passer pour voir~; prévenir bloque au risque de couper. »%
}{Principe du cours}

\niveauChapitre{Intermédiaire}{Chapitre~\ref{chap:siem}~; notions de réseaux.}

\begin{boxobjectifs}
\begin{itemize}
    \item Distinguer IDS et IPS, réseau et hôte.
    \item Comparer détection par signatures et par anomalies.
    \item Comprendre le placement des sondes.
    \item Relier les sondes au SIEM.
    \item Peser le risque d'un blocage automatique.
\end{itemize}
\end{boxobjectifs}

\bigskip

% ============================================================================
\section*{Question directrice : fallait-il bloquer automatiquement, au risque de couper la production de Garoua~?}
% ============================================================================

Un IPS aurait pu bloquer la connexion suspecte du 23~février. Mais un blocage
mal réglé coupe aussi des flux légitimes. Ce chapitre distingue les sondes,
leurs méthodes et leur placement, et pose l'arbitrage entre détecter et
bloquer.

\medskip
\textbf{Plan du chapitre.}
\begin{enumerate}
    \item IDS et IPS, réseau et hôte (section~\ref{sec:c06-types}).
    \item Signatures ou anomalies (section~\ref{sec:c06-methodes}).
    \item Le placement des sondes (section~\ref{sec:c06-placement}).
    \item Détecter ou bloquer (section~\ref{sec:c06-arbitrage}).
\end{enumerate}

\separateur

% ============================================================================
\section{IDS et IPS, Réseau et Hôte}
\label{sec:c06-types}
% ============================================================================

\begin{tableaucours}{L{2.4cm}L{4.8cm}L{5.0cm}}{%
    \tableauHeader{Sonde} &
    \tableauHeader{Agit} &
    \tableauHeader{Observe}
}
IDS & Détecte et alerte, sans bloquer & Le trafic ou l'activité \\
IPS & Détecte et bloque en ligne & Le trafic ou l'activité \\
Réseau (NIDS/NIPS) & Sur un point de passage du réseau & Les flux \\
Hôte (HIDS/HIPS) & Sur une machine & Les processus, fichiers, journaux \\
\end{tableaucours}
\vspace{-0.8em}
\begin{center}{\small\itshape Tableau~6.1. Familles de sondes de détection}\end{center}

% ============================================================================
\section{Signatures ou Anomalies}
\label{sec:c06-methodes}
% ============================================================================

\begin{tableaucours}{L{3.2cm}L{4.6cm}L{4.4cm}}{%
    \tableauHeader{Méthode} &
    \tableauHeader{Principe} &
    \tableauHeader{Limite}
}
Signatures & Reconnaît des motifs connus & Aveugle à l'inconnu \\
Anomalies & Signale l'écart à un profil normal & Faux positifs, apprentissage délicat \\
\end{tableaucours}
\vspace{-0.8em}
\begin{center}{\small\itshape Tableau~6.2. Deux méthodes de détection}\end{center}

Les deux sont complémentaires~: les signatures attrapent le connu à faible
bruit, les anomalies ouvrent une chance sur l'inconnu, au prix de faux
positifs à régler.

% ============================================================================
\section{Le Placement des Sondes}
\label{sec:c06-placement}
% ============================================================================

Une sonde réseau ne voit que le trafic qui passe par elle~: on la place aux
points stratégiques (sortie Internet, cœur de réseau, segment serveurs). Une
sonde hôte voit tout de sa machine mais rien des autres. Les alertes des
sondes alimentent le SIEM, qui les corrèle avec les journaux
(Chapitre~\ref{chap:siem}).

% ============================================================================
\section{Détecter ou Bloquer}
\label{sec:c06-arbitrage}
% ============================================================================

\begin{boxalert}
\textbf{Erreur fréquente.} Mettre tous les IPS en blocage automatique « pour
être tranquille ».

\smallskip
\textbf{Règle~:} un blocage mal réglé coupe des flux légitimes et peut arrêter
la production. On bloque automatiquement le connu à très forte confiance~; on
détecte et on fait valider le reste par un analyste.
\end{boxalert}

\begin{boxscenario}{Opération BÉNOUÉ. Détecter d'abord}
\textbf{Constatation.} LiZbiZ installe une sonde réseau en sortie et des
sondes hôte sur les serveurs. Les sondes sont d'abord en détection, non en
blocage.

\medskip
\textbf{Analyse} (confiance moyenne). Sur un parc sans historique de réglage,
commencer en détection évite de couper la production~; le passage au blocage
se fait règle par règle, une fois le bruit maîtrisé. Prudence d'abord, au nom
de la disponibilité.
\end{boxscenario}

\begin{boxmetier}{Analyste SOC et administrateur réseau}
\textbf{Régler avant de bloquer.} Le passage d'un IPS en blocage se décide
sur mesure~: on observe en détection, on mesure les faux positifs, puis on
bloque les règles sûres. La disponibilité de la production est un critère de
sécurité.
\end{boxmetier}

\separateur

\begin{boxresume}
\textbf{Ce qu'il faut retenir de ce chapitre~:}
\begin{itemize}
    \item \textbf{IDS} alerte, \textbf{IPS} bloque~; \textbf{réseau} ou
          \textbf{hôte}.
    \item \textbf{Signatures} (le connu) et \textbf{anomalies} (l'inconnu)
          se complètent.
    \item \textbf{Le placement} décide de ce qu'une sonde voit.
    \item \textbf{Les sondes alimentent le SIEM}.
    \item \textbf{Bloquer} se fait avec prudence, règle par règle.
\end{itemize}
\end{boxresume}

\begin{boxinfo}
\textbf{Points de contrôle de l'élève.} Je sais~: distinguer IDS/IPS et
réseau/hôte~; comparer signatures et anomalies~; placer une sonde~; arbitrer
entre détecter et bloquer.
\end{boxinfo}

\bigskip

\begin{boxexo}[title={\ding{45}\quad Exercice~6.1. IDS ou IPS \niveaubase}]
Lequel bloque en ligne~? Lequel se contente d'alerter~?
\end{boxexo}

\begin{boxexo}[title={\ding{45}\quad Exercice~6.2. Réseau ou hôte \niveaubase}]
Quelle sonde voit les processus d'une machine~? Laquelle voit les flux~?
\end{boxexo}

\begin{boxexo}[title={\ding{45}\quad Exercice~6.3. Méthodes \niveaubase}]
Quelle méthode est aveugle à l'inconnu~? Laquelle produit des faux positifs~?
\end{boxexo}

\begin{boxexo}[title={\ding{45}\quad Exercice~6.4. Placement \niveaumoyen}]
Où placer une sonde réseau pour voir l'exfiltration vers l'Internet~?
\end{boxexo}

\begin{boxexo}[title={\ding{45}\quad Exercice~6.5. Complémentarité \niveaumoyen}]
Pourquoi combiner signatures et anomalies~?
\end{boxexo}

\begin{boxexo}[title={\ding{45}\quad Exercice~6.6. Blocage \niveaumoyen}]
Pourquoi ne pas tout mettre en blocage automatique chez LiZbiZ~?
\end{boxexo}

\begin{boxexo}[title={\ding{45}\quad Exercice~6.7. SIEM \niveaumoyen}]
Comment les alertes des sondes s'articulent-elles avec le SIEM~?
\end{boxexo}

\begin{boxexo}[title={\ding{45}\quad Exercice~6.8. Angle mort \niveaumoyen}]
Quel trafic une sonde réseau mal placée laisse-t-elle passer~?
\end{boxexo}

\begin{boxexo}[title={\ding{45}\quad Exercice~6.9. Plan de déploiement \niveauexpert}]
Proposez le plan de déploiement des sondes pour LiZbiZ (points, détection
puis blocage).
\end{boxexo}

\begin{boxexo}[title={\ding{45}\quad Exercice~6.10. Passage au blocage \niveauexpert}]
Décrivez la démarche pour faire passer une règle d'IPS de la détection au
blocage sans risque pour la production.
\end{boxexo}

\bigskip

\begin{boxinfo}
\textbf{Pour aller plus loin.}
\begin{itemize}
    \item NIST SP~800-61r2~\cite{nist80061r2}~; MITRE ATT\&CK~\cite{mitreattack}.
\end{itemize}
\end{boxinfo}

\begin{boxinfo}
\textbf{Transition.} Sondes et SIEM ne valent que par leurs règles~: le
Chapitre~\ref{chap:regles} apprend à en écrire et à les éprouver.
\end{boxinfo}
